\newpage
\section*{ДОДАТКИ}
\addcontentsline{toc}{section}{ДОДАТКИ}

\subsection*{Додаток А. Структура репозиторію програмного продукту}
\addcontentsline{toc}{subsection}{Додаток А. Структура репозиторію програмного продукту}

Вихідний код програмного продукту організовано за модульним принципом. Основні каталоги репозиторію:

\begin{itemize}
  \item \texttt{src/core} --- доменні моделі та бізнес-логіка чатів;
  \item \texttt{src/web} --- Flask routes, HTML-шаблони, CSS і JavaScript web-інтерфейсу;
  \item \texttt{src/cli} --- консольний режим роботи;
  \item \texttt{src/storage} --- абстракція сховища та PostgreSQL-реалізація;
  \item \texttt{src/llm} --- LLM-сервіси, runtime-керування моделлю та model registry;
  \item \texttt{migrations} --- Alembic-міграції бази даних;
  \item \texttt{tests} --- автоматизовані тести;
  \item \texttt{scripts} --- допоміжні скрипти для завантаження моделей та адаптерів.
\end{itemize}

Така структура відповідає обраному багаторівневому підходу і дозволяє окремо розвивати інтерфейс користувача, бізнес-логіку, сховище даних та LLM-рівень.

\subsection*{Додаток Б. Фрагменти модульних тестів}
\addcontentsline{toc}{subsection}{Додаток Б. Фрагменти модульних тестів}

Нижче наведено приклади тестів, які перевіряють основні компоненти системи.

\textbf{Тест створення чату у \texttt{ChatManager}:}

\begin{lstlisting}[style=plantumlstyle]
def test_create_chat_with_custom_title(tmp_path):
    manager = make_manager(tmp_path)

    chat = manager.create_chat("Course Project")

    assert chat.title == "Course Project"
    assert manager.get_chat(chat.id).title == "Course Project"
\end{lstlisting}

Фрагмент перевіряє, що доменний сервіс створює чат із заданою назвою та дозволяє отримати його зі сховища.

\textbf{Тест валідації дубльованих назв чатів:}

\begin{lstlisting}[style=plantumlstyle]
def test_reject_duplicate_title_case_insensitively(tmp_path):
    manager = make_manager(tmp_path)
    manager.create_chat("Project Chat")

    with pytest.raises(DuplicateChatTitleError):
        manager.create_chat("project chat")
\end{lstlisting}

Фрагмент перевіряє, що система не дозволяє створювати дублікати назв чатів незалежно від регістру символів.

\textbf{Тест перемикання LoRA/QLoRA-адаптера у runtime:}

\begin{lstlisting}[style=plantumlstyle]
def test_runtime_switch_applies_adapter_path(monkeypatch):
    captured = {}
    manager = make_manager()
    runtime = LLMRuntime(manager)

    def fake_create(backend, model_name_or_path=None,
                    generation_preset=None, adapter_path=None):
        captured["adapter_path"] = adapter_path
        return FakeLLMService()

    monkeypatch.setattr("src.llm.runtime.create_llm_service", fake_create)

    runtime.switch(
        "transformers",
        model_name_or_path="models/qwen",
        adapter_path="adapters/qwen-lora",
    )

    assert captured["adapter_path"] == "adapters/qwen-lora"
\end{lstlisting}

Фрагмент перевіряє, що шлях до adapter передається у фабрику LLM-сервісу під час runtime-перемикання моделі.

\textbf{Тест web endpoint для надсилання повідомлення:}

\begin{lstlisting}[style=plantumlstyle]
def test_send_endpoint_adds_user_and_assistant_messages(tmp_path):
    client, manager = make_client(tmp_path)
    chat = manager.create_chat("Web chat")

    response = client.post(f"/chat/{chat.id}/send",
                           json={"message": "Hello"})

    payload = response.get_json()
    assert response.status_code == 200
    assert payload["ok"] is True
    assert payload["user_message"]["content"] == "Hello"
    assert payload["assistant_message"]["role"] == "assistant"
    assert len(manager.get_chat(chat.id).messages) == 2
\end{lstlisting}

Фрагмент перевіряє, що HTTP endpoint додає повідомлення користувача, генерує відповідь асистента і зберігає обидва повідомлення в історії чату.

\textbf{Тест сканування локальної моделі:}

\begin{lstlisting}[style=plantumlstyle]
def test_folder_with_config_and_safetensors_is_listed(tmp_path):
    model_dir = tmp_path / "models" / "qwen"
    model_dir.mkdir(parents=True)
    (model_dir / "config.json").write_text("{}", encoding="utf-8")
    (model_dir / "model.safetensors").write_text("fake", encoding="utf-8")
    (model_dir / "tokenizer.json").write_text("{}", encoding="utf-8")

    models = list_local_models(tmp_path / "models")

    assert len(models) == 1
    assert models[0].name == "qwen"
    assert models[0].has_config is True
    assert models[0].has_weights is True
    assert models[0].has_tokenizer is True
\end{lstlisting}

Фрагмент перевіряє, що model registry знаходить локальну модель лише тоді, коли в каталозі є конфігурація, ваги та tokenizer.

\subsection*{Додаток В. Результати покриття коду тестами}
\addcontentsline{toc}{subsection}{Додаток В. Результати покриття коду тестами}

Для оцінювання покриття було використано \texttt{coverage.py}. Повний запуск тестів:

\begin{lstlisting}[style=plantumlstyle]
python -m coverage run -m pytest
\end{lstlisting}

Результат виконання:

\begin{lstlisting}[style=plantumlstyle]
105 passed in 3.26s
\end{lstlisting}

Звіт покриття для production-коду було сформовано командою:

\begin{lstlisting}[style=plantumlstyle]
python -m coverage report --include="src/*,migrations/*,scripts/*"
\end{lstlisting}

Основні результати:

\begin{longtable}{|p{8cm}|p{3cm}|}
\caption*{Фрагмент звіту покриття production-коду}\\
\hline
\textbf{Модуль} & \textbf{Покриття} \\
\hline
\texttt{src/core/chat\_manager.py} & 94\% \\
\hline
\texttt{src/config.py} & 98\% \\
\hline
\texttt{src/llm/runtime.py} & 97\% \\
\hline
\texttt{src/llm/model\_registry.py} & 92\% \\
\hline
\texttt{src/storage/postgres\_storage.py} & 77\% \\
\hline
\texttt{src/web/web\_app.py} & 81\% \\
\hline
\textbf{Загальне покриття production-коду} & \textbf{77\%} \\
\hline
\end{longtable}

Отримане покриття відповідає рекомендованому рівню 70--80\% і підтверджує, що основна бізнес-логіка та ключові інтеграційні точки системи перевіряються автоматизованими тестами.

\subsection*{Додаток Г. Приклади функціонального тестування web/API}
\addcontentsline{toc}{subsection}{Додаток Г. Приклади функціонального тестування web/API}

Функціональні сценарії перевірялися через Flask test client. Вони імітують HTTP-запити користувача до web-застосунку.

\begin{longtable}{|p{3.5cm}|p{2cm}|p{4.8cm}|p{4cm}|}
\caption*{Приклади функціонального тестування}\\
\hline
\textbf{Дія} & \textbf{Метод} & \textbf{Endpoint} & \textbf{Очікуваний результат} \\
\hline
Відкрити головну сторінку & GET & \texttt{/} & Статус 200, відображення sidebar, model status і списку чатів. \\
\hline
Створити чат & POST & \texttt{/chats} & Новий чат створено, користувача перенаправлено на сторінку чату. \\
\hline
Надіслати повідомлення & POST & \texttt{/chat/<id>/send} & Збережено повідомлення користувача й відповідь асистента. \\
\hline
Streaming-генерація & POST & \texttt{/chat/<id>/stream} & Відповідь повертається NDJSON-частинами та зберігається після завершення. \\
\hline
Зупинити генерацію & POST & \texttt{/api/generation/stop} & Активний stop event встановлено, генерація завершується контрольовано. \\
\hline
Перемкнути модель & POST & \texttt{/api/model/switch} & Runtime переходить у стан завантаження, після успіху активна модель оновлюється. \\
\hline
Пошук повідомлень & GET & \texttt{/api/messages/search} & Повертаються релевантні повідомлення з назвою чату і роллю автора. \\
\hline
Експорт чату & GET & \texttt{/chat/<id>/export} & Користувач отримує Markdown-файл з історією діалогу. \\
\hline
\end{longtable}

Приклад перевірки помилки валідації порожнього повідомлення:

\begin{lstlisting}[style=plantumlstyle]
def test_send_endpoint_rejects_empty_message(tmp_path):
    client, manager = make_client(tmp_path)
    chat = manager.create_chat("Web chat")

    response = client.post(f"/chat/{chat.id}/send",
                           json={"message": "   "})

    assert response.status_code == 400
    assert response.get_json() == {
        "ok": False,
        "error": "Message cannot be empty.",
    }
\end{lstlisting}

Цей тест перевіряє, що web-рівень не приймає порожні повідомлення і повертає користувачу коректну помилку.

